\documentclass[12pt,a4paper]{article}

\usepackage[T2A]{fontenc}
\usepackage[utf8]{inputenc}
\usepackage[russian]{babel}
\usepackage{geometry}
\usepackage{setspace}
\usepackage{fancyvrb}
\usepackage{titlesec}
\usepackage{hyperref}

\geometry{left=25mm,right=15mm,top=20mm,bottom=20mm}
\setstretch{1.2}
\titleformat{\section}{\large\bfseries}{\thesection.}{0.5em}{}

\begin{document}

\begin{titlepage}
\begin{center}
\textbf{Федеральное государственное образовательное учреждение высшего образования}\\
\textbf{<<Национальный исследовательский университет ИТМО>>}\\[1.5cm]

Факультет программной инженерии и компьютерной техники\\
Дисциплина <<ОПИ>>\\[2.5cm]

\textbf{\Large Отчёт по лабораторной работе №3}\\[0.3cm]
\textbf{\Large Сценарий сборки Apache Ant}\\[0.8cm]
\normalsize
Тема: компиляция, тестирование и упаковка проекта в jar-архив\\
с дополнительными целями автоматизации
\end{center}

\vfill

\begin{flushright}
\textbf{Выполнил:}\\
Смирнов Вадим Константинович\\
Группа: P3119\\[0.7cm]

\textbf{Проверил: Ермаков М.К.}\\
\underline{\hspace{6cm}}
\end{flushright}

\vfill

\begin{center}
г. Санкт-Петербург\\
2026
\end{center}
\end{titlepage}

\tableofcontents
\newpage

\section{Цель и постановка работы}
Цель лабораторной работы --- освоить средства мониторинга и профилирования JVM-приложения
на примере web-приложения из лабораторной работы №3 по дисциплине <<Веб-программирование>>.

В рамках работы требовалось:
\begin{itemize}
    \item реализовать два MBean-класса для пользовательских метрик;
    \item выполнить мониторинг приложения в JConsole;
    \item выполнить мониторинг и анализ памяти в VisualVM;
    \item локализовать и устранить проблему производительности и зафиксировать алгоритм действий.
\end{itemize}

\section{Реализация MBean-классов}

\subsection{MBean для подсчета точек и оповещений}
Реализован MBean \texttt{PointCountersMBean} с атрибутами:
\begin{itemize}
    \item \texttt{TotalPoints} --- общее число сохраненных точек;
    \item \texttt{HitPoints} --- число попаданий в область;
    \item \texttt{OutOfBoundsEvents} --- число выходов координат за отображаемую область.
\end{itemize}

Класс \texttt{PointCounters} наследует \texttt{NotificationBroadcasterSupport} и отправляет JMX-оповещение
типа \texttt{org.example.point.outOfBounds} при попытке поставить точку вне диапазона
отображаемой плоскости.

\begin{figure}[h]
\centering
\fbox{\parbox[c][6.2cm][c]{0.9\textwidth}{\centering
Скриншот 1 --- код/структура \texttt{PointCountersMBean} и \texttt{PointCounters}\\
(вставить изображение)
}}
\caption{Реализация MBean для счетчиков точек}
\end{figure}

\subsection{MBean для процента попаданий}
Реализован MBean \texttt{HitRatioMBean} с атрибутами:
\begin{itemize}
    \item \texttt{HitPercent} --- процент попаданий;
    \item \texttt{TotalClicks} --- общее число кликов по координатной плоскости.
\end{itemize}

Формула расчета:
\[
HitPercent = \frac{HitPoints}{TotalClicks} \cdot 100\%
\]

\begin{figure}[h]
\centering
\fbox{\parbox[c][6.2cm][c]{0.9\textwidth}{\centering
Скриншот 2 --- код \texttt{HitRatioMBean} и \texttt{HitRatio}\\
(вставить изображение)
}}
\caption{Реализация MBean для отношения попаданий}
\end{figure}

\subsection{Интеграция в приложение}
Для регистрации MBean в \texttt{PlatformMBeanServer} используется класс \texttt{MBeanRegistrar}.
Регистрация выполняется по именам:
\begin{itemize}
    \item \texttt{org.example.monitoring:type=PointCounters};
    \item \texttt{org.example.monitoring:type=HitRatio}.
\end{itemize}

Сервис \texttt{PointMetricsService} хранит счетчики (\texttt{AtomicLong}) и инкапсулирует бизнес-метрики.
В \texttt{PointController} при обработке запроса фиксируются:
\begin{itemize}
    \item клик по графику;
    \item сохранение точки и признак попадания;
    \item событие выхода координат за границы отображаемой области.
\end{itemize}

\begin{figure}[h]
\centering
\fbox{\parbox[c][6.2cm][c]{0.9\textwidth}{\centering
Скриншот 3 --- регистрация MBean в \texttt{MBeanRegistrar} и вызовы в \texttt{PointController}\\
(вставить изображение)
}}
\caption{Подключение метрик к пользовательскому сценарию}
\end{figure}

\section{Мониторинг в JConsole}

\subsection{Ход выполнения}
\begin{enumerate}
    \item Запущено web-приложение и подключена утилита JConsole к JVM-процессу.
    \item Вкладка \texttt{MBeans}: проверено наличие \texttt{PointCounters} и \texttt{HitRatio}.
    \item Выполнены действия в приложении (клики, попадания и выходы за область).
    \item Сняты значения атрибутов MBean.
    \item Во вкладке \texttt{Runtime} зафиксировано время с момента запуска JVM.
\end{enumerate}

\begin{figure}[h]
\centering
\fbox{\parbox[c][6.2cm][c]{0.9\textwidth}{\centering
Скриншот 4 --- подключение JConsole к процессу приложения\\
(вставить изображение)
}}
\caption{Подключение JConsole}
\end{figure}

\begin{figure}[h]
\centering
\fbox{\parbox[c][6.2cm][c]{0.9\textwidth}{\centering
Скриншот 5 --- значения атрибутов \texttt{PointCounters} и \texttt{HitRatio} в JConsole\\
(вставить изображение)
}}
\caption{Показания MBean в JConsole}
\end{figure}

\begin{figure}[h]
\centering
\fbox{\parbox[c][6.2cm][c]{0.9\textwidth}{\centering
Скриншот 6 --- вкладка Runtime (Uptime в миллисекундах)\\
(вставить изображение)
}}
\caption{Время работы JVM}
\end{figure}

\subsection{Зафиксированные значения}
\begin{itemize}
    \item \texttt{TotalClicks}: \underline{\hspace{5cm}}
    \item \texttt{TotalPoints}: \underline{\hspace{5cm}}
    \item \texttt{HitPoints}: \underline{\hspace{5cm}}
    \item \texttt{OutOfBoundsEvents}: \underline{\hspace{5cm}}
    \item \texttt{HitPercent}: \underline{\hspace{5cm}}
    \item \texttt{Uptime, ms}: \underline{\hspace{5cm}}
\end{itemize}

\section{Мониторинг и профилирование в VisualVM}

\subsection{Графики изменения MBean-метрик}
В VisualVM выполнено подключение к приложению и добавлен плагин JMX (при необходимости).
Во время работы с интерфейсом приложения снят график изменения метрик:
\texttt{TotalClicks}, \texttt{TotalPoints}, \texttt{HitPoints}, \texttt{OutOfBoundsEvents}, \texttt{HitPercent}.

\begin{figure}[h]
\centering
\fbox{\parbox[c][6.2cm][c]{0.9\textwidth}{\centering
Скриншот 7 --- график изменения MBean-показаний во времени в VisualVM\\
(вставить изображение)
}}
\caption{Динамика пользовательских метрик}
\end{figure}

\subsection{Анализ памяти}
С помощью профайлера памяти VisualVM получен heap histogram и определены:
\begin{itemize}
    \item класс, объекты которого занимают наибольший объем памяти;
    \item пользовательский класс, чьи экземпляры удерживают данные объекты.
\end{itemize}

\begin{figure}[h]
\centering
\fbox{\parbox[c][6.2cm][c]{0.9\textwidth}{\centering
Скриншот 8 --- heap histogram (класс с максимальным объемом)\\
(вставить изображение)
}}
\caption{Результат анализа распределения памяти}
\end{figure}

\begin{figure}[h]
\centering
\fbox{\parbox[c][6.2cm][c]{0.9\textwidth}{\centering
Скриншот 9 --- путь удержания/ссылка на пользовательский класс\\
(вставить изображение)
}}
\caption{Связь с пользовательским классом}
\end{figure}

\subsection{Зафиксированные результаты}
\begin{itemize}
    \item Класс с наибольшим объемом памяти: \underline{\hspace{7cm}}
    \item Пользовательский класс-владелец: \underline{\hspace{7cm}}
    \item Краткий комментарий: \underline{\hspace{7cm}}
\end{itemize}

\section{Локализация и устранение проблемы производительности}

\subsection{Описание выявленной проблемы}
В ходе профилирования выявлена проблема производительности:
\underline{\hspace{12cm}}
\newline
\underline{\hspace{12cm}}
\newline
\underline{\hspace{12cm}}

\subsection{Способ устранения}
Для устранения проблемы были выполнены следующие действия:
\begin{enumerate}
    \item \underline{\hspace{11cm}}
    \item \underline{\hspace{11cm}}
    \item \underline{\hspace{11cm}}
\end{enumerate}

\begin{figure}[h]
\centering
\fbox{\parbox[c][6.2cm][c]{0.9\textwidth}{\centering
Скриншот 10 --- профилирование до исправления (узкое место)\\
(вставить изображение)
}}
\caption{Состояние до оптимизации}
\end{figure}

\begin{figure}[h]
\centering
\fbox{\parbox[c][6.2cm][c]{0.9\textwidth}{\centering
Скриншот 11 --- изменения в коде, устраняющие проблему\\
(вставить изображение)
}}
\caption{Внесенные исправления}
\end{figure}

\begin{figure}[h]
\centering
\fbox{\parbox[c][6.2cm][c]{0.9\textwidth}{\centering
Скриншот 12 --- профилирование после исправления\\
(вставить изображение)
}}
\caption{Состояние после оптимизации}
\end{figure}

\subsection{Алгоритм локализации проблемы}
\begin{enumerate}
    \item Запуск приложения и подключение профайлера VisualVM/IDE.
    \item Воспроизведение пользовательского сценария с характерной нагрузкой.
    \item Снятие CPU и heap snapshot.
    \item Выделение методов/классов с наибольшими затратами времени или памяти.
    \item Проверка причин (частота вызовов, аллокации, лишние операции).
    \item Внесение изменений в код и повторный замер на тех же действиях.
    \item Сравнение результатов <<до/после>> и фиксация улучшений.
\end{enumerate}

\section{Вывод}
В ходе лабораторной работы реализованы и подключены два MBean-класса для пользовательских метрик,
проведен мониторинг в JConsole и VisualVM, а также выполнены действия по выявлению и устранению
проблемы производительности. Полученные результаты подтверждают корректность работы JMX-мониторинга
и применимость инструментов профилирования для анализа web-приложений на JVM.

\end{document}
